\section{Analysis}
\label{sec:analysis}

This section presents four complementary analyses: statistical significance, data efficiency, computational cost, and a synthesis of practical guidelines.

\subsection{Statistical Significance}
\label{sec:significance}

To confirm that the observed accuracy differences are not due to random variation, we run each method five times with different random seeds and conduct paired $t$-tests and Wilcoxon signed-rank tests comparing each method against LDA.

Table~\ref{tab:significance} summarizes selected comparisons (for each backbone--dataset pair, we show the three most informative contrasts). LDA's advantage over full features is significant at $p < 0.001$ in all eight configurations. The gain over PCA is significant in seven of eight ($p = 0.27$ for ResNet-18/Tiny ImageNet, where the two methods perform nearly identically). LFDA is significantly \emph{worse} than LDA everywhere (e.g., $\Delta = -4.04\%$ on ResNet-50/CIFAR-100). Among our proposed extensions, DSB significantly outperforms LDA in six of eight cases ($+0.30$ to $+0.41\%$), but on EfficientNet/Tiny ImageNet it is \emph{significantly worse} ($\Delta = -0.20\%$, $p = 0.002$), indicating that iterative scatter reweighting does not uniformly help.

\begin{table}[t]
\centering
\caption{Multi-seed statistical significance (5 seeds). Each row compares a method against LDA using a paired $t$-test. $\Delta$: mean accuracy difference (method $-$ LDA). Bold $p$-values indicate statistical significance at $\alpha = 0.05$. We show all six comparison methods; the pattern is consistent across all four backbones (see supplementary for full table).}
\label{tab:significance}
\small
\setlength{\tabcolsep}{3pt}
\begin{tabular}{@{}l l l r r@{}}
\toprule
Backbone & Dataset & Method & $\Delta$(\%) & $p$-value \\
\midrule
\multirow{6}{*}{ResNet-18}
& \multirow{3}{*}{C100}  & Full  & $-$4.12 & $\mathbf{<.001}$ \\
&                         & PCA   & $-$1.91 & $\mathbf{<.001}$ \\
&                         & DSB   & $+$0.30 & $\mathbf{<.001}$ \\
& \multirow{3}{*}{TinyIN} & Full  & $-$4.57 & $\mathbf{<.001}$ \\
&                         & PCA   & $+$0.04 & .273           \\
&                         & R-LDA & $+$0.32 & $\mathbf{<.001}$ \\
\midrule
\multirow{6}{*}{ResNet-50}
& \multirow{3}{*}{C100}  & Full  & $-$0.26 & $\mathbf{<.001}$ \\
&                         & LFDA  & $-$4.04 & $\mathbf{<.001}$ \\
&                         & RDA   & $+$0.43 & $\mathbf{<.001}$ \\
& \multirow{3}{*}{TinyIN} & Full  & $-$0.72 & $\mathbf{<.001}$ \\
&                         & PCA   & $-$1.66 & $\mathbf{<.001}$ \\
&                         & RDA   & $+$0.03 & .303           \\
\midrule
\multirow{6}{*}{MobNetV3}
& \multirow{3}{*}{C100}  & Full  & $-$2.89 & $\mathbf{<.001}$ \\
&                         & PCA   & $-$3.86 & $\mathbf{<.001}$ \\
&                         & DSB   & $+$0.41 & $\mathbf{<.001}$ \\
& \multirow{3}{*}{TinyIN} & Full  & $-$4.11 & $\mathbf{<.001}$ \\
&                         & PCA   & $-$1.40 & $\mathbf{<.001}$ \\
&                         & DSB   & $+$0.38 & $\mathbf{<.001}$ \\
\midrule
\multirow{6}{*}{EffNet}
& \multirow{3}{*}{C100}  & Full  & $-$0.89 & $\mathbf{<.001}$ \\
&                         & PCA   & $-$2.53 & $\mathbf{<.001}$ \\
&                         & RDA   & $+$0.05 & .075           \\
& \multirow{3}{*}{TinyIN} & Full  & $-$0.55 & $\mathbf{<.001}$ \\
&                         & PCA   & $-$0.35 & $\mathbf{.015}$  \\
&                         & DSB   & $-$0.20 & $\mathbf{.002}$  \\
\bottomrule
\end{tabular}
\end{table}

Three configurations show \emph{non-significant} differences between RDA and LDA (ResNet-18/Tiny ImageNet: $p = 0.058$; ResNet-50/Tiny ImageNet: $p = 0.303$; EfficientNet/CIFAR-100: $p = 0.075$). This indicates that RDA's residual augmentation provides inconsistent benefits---the extra PCA dimensions capture noise as often as signal. In the significant cases, RDA sometimes helps (ResNet-50/CIFAR-100: $+0.43$) and sometimes hurts (EfficientNet/Tiny ImageNet: $-0.47$), reinforcing the conclusion that plain LDA is the more reliable choice.

\textit{Note on variance:} Several methods exhibit near-zero standard deviations across seeds (e.g., LDA on ResNet-18/CIFAR-100: all five seeds yield 66.97\%). This occurs because LDA's projection is deterministic and the logistic regression solver converges to the same optimum regardless of initialization. While this makes $t$-statistics numerically large, it reflects genuine reproducibility rather than a statistical artifact. We additionally confirm all key comparisons with the non-parametric Wilcoxon signed-rank test, which makes no distributional assumptions.

\subsection{Data Efficiency}
\label{sec:data_efficiency}

A natural question is how much labeled data LDA needs to outperform the full-feature baseline. We subsample the training set at fractions of 10\%, 25\%, 50\%, and 100\%, repeating each configuration three times with different random subsamples.

\begin{table}[t]
\centering
\caption{Accuracy (\%) at different training set fractions on CIFAR-100 (3-seed mean). \textbf{Bold}: best method per fraction.}
\label{tab:data_efficiency}
\small
\begin{tabular}{@{}l l rrrr@{}}
\toprule
Backbone & Method & 10\% & 25\% & 50\% & 100\% \\
\midrule
\multirow{4}{*}{ResNet-18}
& Full & \textbf{57.7} & \textbf{60.3} & 61.6 & 62.9 \\
& PCA  & 51.4 & 58.4 & 62.5 & 65.1 \\
& LDA  & 54.7 & 58.8 & 63.9 & 67.0 \\
& DSB  & 55.3 & 58.4 & \textbf{64.2} & \textbf{67.3} \\
\midrule
\multirow{4}{*}{ResNet-50}
& Full & \textbf{66.8} & \textbf{70.0} & \textbf{71.2} & 72.0 \\
& PCA  & 59.2 & 63.6 & 66.9 & 69.1 \\
& LDA  & 62.9 & 66.9 & 69.4 & 72.2 \\
& DSB  & 64.3 & 67.0 & 69.3 & \textbf{72.6} \\
\bottomrule
\end{tabular}
\end{table}

Table~\ref{tab:data_efficiency} reveals an important pattern. At 10\% of the training data (5,000 images for CIFAR-100), full features outperform all reduction methods. This makes sense: LDA estimates $C \times D$ parameters for the within-class scatter matrix, and with only 50 samples per class, these estimates are noisy. PCA suffers even more because it discards variance directions that happen to be discriminative.

The crossover occurs between 25\% and 50\% of the data. At 50\% (25,000 images), LDA overtakes full features on ResNet-18 by 2.3 points and closes the gap on ResNet-50. At 100\%, LDA's advantage is clear.

This has a practical implication: \emph{if your target dataset has fewer than approximately 50 samples per class, skip dimensionality reduction and classify the full features directly.} Above that threshold, LDA is the better choice.

For ResNet-50, the crossover is later because the 2048-dimensional full features are already quite effective, and the scatter matrix estimation is harder in higher dimensions. This underscores a general principle: the benefit of supervised reduction scales with the gap between the feature dimensionality and the intrinsic task complexity.

Fig.~\ref{fig:data_efficiency} plots the accuracy--fraction curves, clearly showing the crossover point.

\begin{figure}[t]
    \centering
    \includegraphics[width=\columnwidth]{fig4_data_efficiency.pdf}
    \caption{Accuracy vs.\ training set fraction on CIFAR-100 for ResNet-18 (left) and ResNet-50 (right). Full features dominate at 10\%, but LDA overtakes between 25--50\% and widens the gap at 100\%.}
    \label{fig:data_efficiency}
\end{figure}

\subsection{Computational Cost Analysis}
\label{sec:cost}

A method's value depends not just on its accuracy but on the cost of achieving it. We analyze this through \emph{Pareto optimality}: a method is Pareto-optimal if no other method is simultaneously faster \emph{and} more accurate.

\begin{table}[t]
\centering
\caption{Average accuracy and time across all 8 backbone--dataset configurations. ``Pareto'' counts configurations (out of 8) where the method is Pareto-optimal.}
\label{tab:efficiency}
\small
\begin{tabular}{@{}l rrr@{}}
\toprule
Method & Avg. Acc. (\%) & Avg. Time (s) & Pareto (/8) \\
\midrule
PCA    & 67.5 & 27 & 8 \\
\textbf{LDA}    & \textbf{69.4} & \textbf{35} & \textbf{7} \\
R-LDA  & 69.2 & 46 & 3 \\
LFDA   & 67.5 & 39 & 0 \\
RDA    & 69.4 & 49 & 6 \\
DSB    & 69.6 & 87 & 4 \\
Full   & 67.1 & 81 & 0 \\
\bottomrule
\end{tabular}
\end{table}

Table~\ref{tab:efficiency} reveals the accuracy--cost landscape. PCA is Pareto-optimal in all eight configurations by virtue of its speed, but trails LDA by nearly 2 points in accuracy. LDA is Pareto-optimal in seven of eight---it achieves the same accuracy as RDA while being 1.4$\times$ faster, and is within 0.2\% of DSB at 2.5$\times$ lower cost.

Crucially, \textbf{full features are never Pareto-optimal}: they are both slower (due to higher-dimensional logistic regression) and less accurate than LDA in every configuration. This counterintuitive result---that discarding 80--95\% of dimensions \emph{improves} both speed and accuracy---is the central practical finding of this paper.

LFDA is also never Pareto-optimal: it is slower than LDA while being less accurate. NCA is dominated by every other method due to its extreme training time (260--330s). Among extensions, DSB is Pareto-optimal in four configurations where its accuracy gain justifies the 2--3$\times$ time cost.

Fig.~\ref{fig:pareto} visualizes the accuracy--time tradeoff, with the Pareto frontier shown as a dashed line. LDA sits close to the frontier on both datasets.

\begin{figure}[t]
    \centering
    \includegraphics[width=\columnwidth]{fig2_pareto.pdf}
    \caption{Accuracy vs.\ wall-clock time (averaged across four backbones) for CIFAR-100 (left) and Tiny ImageNet (right). Dashed line: Pareto frontier. Full features are dominated---slower \emph{and} less accurate than LDA.}
    \label{fig:pareto}
\end{figure}

\subsection{Effect of Number of Components}
\label{sec:components}

Our main experiments use $d = C-1$ components, the theoretical maximum for LDA. To validate this choice, we sweep $d \in \{5, 10, 20, 40, 60, 80, 99\}$ for both LDA and PCA on CIFAR-100 with two backbones (3 seeds each).

Fig.~\ref{fig:component_sweep} shows that accuracy increases monotonically with $d$ for both methods, with LDA dominating PCA at every dimensionality. On ResNet-18, LDA at $d=99$ achieves 66.97\%, surpassing the full-feature baseline (62.85\%) by over 4 points. On ResNet-50, LDA at $d=99$ reaches 72.19\%, matching the full baseline (72.05\%). The curves show no sign of overfitting or diminishing returns as $d$ approaches $C-1$, confirming that \textbf{using the maximum number of LDA components is optimal}.

The LDA--PCA gap widens with $d$: at $d=5$, LDA leads by 2--7 points; at $d=99$, the gap is 1.9--4.1 points. This suggests that the later LDA components (higher-order discriminant directions) carry less discriminative information but still contribute positively.

\begin{figure}[t]
    \centering
    \includegraphics[width=\columnwidth]{fig3_component_sweep.pdf}
    \caption{Accuracy vs.\ number of projected dimensions ($d$) on CIFAR-100 for LDA and PCA, with full-feature baselines (dashed). Both LDA and PCA improve monotonically; $d = C-1 = 99$ is optimal.}
    \label{fig:component_sweep}
\end{figure}

\subsection{Why Does LDA Help?}

The consistent benefit of LDA over full features warrants explanation. We identify three complementary mechanisms:

\paragraph{Implicit regularization.} By projecting from $D$ to $C-1$ dimensions, LDA eliminates feature dimensions that contribute to overfitting in the downstream classifier. The logistic regression in $D = 2048$ dimensions with $\ell_2$ regularization still has more capacity to overfit than in $d = 99$ dimensions. LDA acts as a structured alternative to stronger $\ell_2$ penalties, one that uses label information to select which dimensions to keep.

\paragraph{Removal of task-irrelevant variance.} ImageNet features encode information about 1,000 categories. For a 100-class target task, much of this information is irrelevant. PCA retains high-variance directions regardless of task relevance; LDA explicitly retains directions that separate the target classes.

\paragraph{Improved conditioning for the classifier.} High-dimensional feature spaces often contain near-collinear feature groups, which slow convergence of iterative solvers like LBFGS. LDA produces orthogonal, decorrelated features that are better conditioned, leading to faster and more stable classifier training.

The magnitude of LDA's gain correlates strongly with the ratio $D / (C-1)$: how much the full feature space exceeds the discriminant subspace. Across our eight configurations, we observe a Pearson correlation of $r = -0.78$ between $D/(C-1)$ and LDA's accuracy gain over Full features. When $D/(C-1)$ is \emph{low}---meaning the backbone is compact relative to the number of classes (e.g., ResNet-18 on Tiny ImageNet, $D/(C-1) = 2.6$)---the full-feature classifier is already stretched thin, and LDA provides the greatest lift (+4.58\%). Conversely, when $D/(C-1)$ is \emph{high} (e.g., ResNet-50 on CIFAR-100, $D/(C-1) = 20.7$), the full-feature space has ample capacity and LDA's gain is modest (+0.26\%).

This suggests that LDA's primary value is not in reducing ``excess'' dimensions from oversized feature spaces, but rather in \emph{concentrating information} when the full-feature classifier is under-performing---precisely the regime where practitioners need the most help.

\subsection{When Does LDA \textit{Not} Help?}

Our results identify two regimes where LDA's advantage is minimal:

\begin{enumerate}[leftmargin=*]
    \item \textbf{Very small training sets} ($< 50$ samples per class). LDA's scatter matrix estimates become unreliable, and full features with strong regularization are preferable (Section~\ref{sec:data_efficiency}).

    \item \textbf{High-capacity backbones on few-class tasks.} When $D/(C-1)$ is large (e.g., ResNet-50 at 2048D on 100-class CIFAR-100, ratio $= 20.7$), the full features already carry enough discriminative capacity that $\ell_2$-regularized logistic regression performs well, and LDA's projection offers only marginal improvement. This explains why PCA can also match LDA in certain configurations---on ResNet-18/Tiny ImageNet ($D/(C-1) = 2.6$), both PCA and LDA achieve $\approx 64.4\%$ because the 2.6$\times$ reduction ratio leaves little room to differentiate supervised from unsupervised projection.
\end{enumerate}

\subsection{Practitioner Guidelines}
\label{sec:guidelines}

Based on our comprehensive evaluation, we offer the following concrete recommendations:

\begin{enumerate}[leftmargin=*]
    \item \textbf{Always try LDA before deploying full features.} In every configuration we tested, LDA either matched or improved accuracy while simultaneously reducing dimensionality and accelerating inference. The risk is minimal: we observed no configuration where LDA degraded performance relative to full features.

    \item \textbf{Set $d = C - 1$.} This is the maximum number of LDA components and consistently the best choice (Section~\ref{sec:components}). Lower values sacrifice accuracy without meaningful speed gains.

    \item \textbf{Skip PCA+LDA.} Two-stage reduction provides no benefit over direct LDA and adds unnecessary complexity. Modern LDA implementations handle the conditioning internally.

    \item \textbf{Skip LFDA and NCA.} Neither improves on LDA for frozen CNN features, and both are substantially slower.

    \item \textbf{Consider DSB only if 0.2--0.4\% accuracy matters} and doubling the reduction time is acceptable. For most applications, the marginal gain does not justify the cost.

    \item \textbf{Use R-LDA when $D \gg C$ and the training set is moderate.} Regularization helps on high-dimensional backbones (ResNet-50, EfficientNet) with Tiny ImageNet's 200 classes.

    \item \textbf{If the training set has $< 50$ samples per class, use full features.} LDA's scatter estimates need sufficient data to be reliable.

    \item \textbf{Standardize after projection.} Feature scaling before the classifier is essential for fair and stable results regardless of the reduction method.
\end{enumerate}
